\section{Complete cluster counts from a finite moment window}
\label{sec:spectral}

The first result isolates a useful form of discrete rigidity. A missed spectral
point contributes a definite positive amount to a carefully chosen residual.
Once the entire residual is smaller than that amount, the missed count must
be zero. A related interpolation construction recovers the integer count in
each proposed interval. The underlying spectrum may be infinite.

\subsection{Setup and the finite data}

Let $(x_n)$ be a finite or countably infinite multiset of positive real numbers,
with repeated values listed according to their positive integer multiplicities,
such that
\begin{equation}\label{eq:summability}
 p_1:=\sum_n x_n<\infty,
 \qquad p_j:=\sum_n x_n^j\quad(j\ge1).
\end{equation}
Every $x_n\le p_1$, so every $p_j$ is finite. There are only finitely many
occurrences above any positive threshold.

Choose $m\ge1$ distinct positive candidate centers $a_1,\ldots,a_m$, positive
radii $r_i<a_i$, and pairwise disjoint open intervals
\[
 I_i=(a_i-r_i,a_i+r_i).
\]
Fix an integer $K\ge1$. Define
\begin{align}
 P(x)&=\prod_{i=1}^m(x-a_i),
 &D_i&=\prod_{j\ne i}(a_i-a_j),\label{eq:P}\\
 L_i(x)&=\frac{\prod_{j\ne i}(x-a_j)}{D_i},
 &d_i&=\frac K{a_i}+2\sum_{j\ne i}\frac1{a_i-a_j},\label{eq:L}\\
 Q_i(x)&=\left(\frac x{a_i}\right)^K
       \bigl[1-d_i(x-a_i)\bigr]L_i(x)^2.\label{eq:Q}
\end{align}
Here and below empty products equal one and empty sums equal zero.
The residual and count estimates are
\begin{equation}\label{eq:RS}
 R_K=\sum_n x_n^K P(x_n)^2\ge0,
 \qquad S_i=\sum_n Q_i(x_n).
\end{equation}
They are absolutely convergent: each summand is a polynomial in $x_n$ with
no constant term, evaluated on a bounded summable spectrum.

The construction has three relevant properties:
\[
 Q_i(a_j)=\begin{cases}1,&j=i,\\0,&j\ne i,\end{cases}
 \qquad Q_i'(a_j)=0,
 \qquad x^K\mid Q_i(x).
\]
At $a_i$, the chosen linear factor cancels the logarithmic derivative
$K/a_i+2L_i'(a_i)$; at other centers, the squared factors give double zeros.
Moreover $\deg Q_i\le K+2m-1$ and $\deg(x^KP^2)=K+2m$. Thus all the
quantities in \eqref{eq:RS} use only the $2m+1$ moments
\begin{equation}\label{eq:window}
 p_K,p_{K+1},\ldots,p_{K+2m}.
\end{equation}
The number of entries in this window is fixed when $m$ is fixed. The starting
order $K$ and the required accuracy need not be small.

\subsection{Explicit count bounds}

For each $i$ set
\begin{align}
 b_i&=a_i-r_i,&g_{ij}&=|a_i-a_j|-r_i\quad(j\ne i),\\
 A_i&=\prod_{j\ne i}g_{ij}^{-2},
 &T_i&=\frac K{b_i}+2\sum_{j\ne i}g_{ij}^{-1},\\
 U_i&=\frac K{b_i^2}+2\sum_{j\ne i}g_{ij}^{-2}.&&
\end{align}
All $b_i,g_{ij}$ are positive by the interval assumptions. Put
\begin{equation}\label{eq:Ci}
 C_i=\max\left\{
 \frac{b_i^{-K}A_i}{2}(T_i^2+U_i),\quad
 \frac{a_i^{-K}}{D_i^2}\left(r_i^{-2}+|d_i|r_i^{-1}\right)
 \right\}.
\end{equation}

\begin{theorem}[Residual certificate]\label{thm:residual}
Under \eqref{eq:summability}, let $N_i$ count all occurrences in $I_i$,
with multiplicity. Then
\begin{equation}\label{eq:countbound}
 |N_i-S_i|\le C_iR_K.
\end{equation}
For any $\tau>0$, let $N_{\rm miss}$ count occurrences with $x_n\ge\tau$
outside $\bigcup_i I_i$. Then
\begin{equation}\label{eq:missing}
 N_{\rm miss}\le\frac{R_K}{\tau^K\prod_i r_i^2}.
\end{equation}
Consequently $R_K<\tau^K\prod_i r_i^2$ excludes every such missed point.
If, in addition to this strict no-missing inequality, the intervals lie
above $\tau$ and each interval
$[S_i-C_iR_K,S_i+C_iR_K]$ contains exactly one nonnegative integer $n_i$,
then the complete spectrum at or above $\tau$ is contained in these intervals
and has the exact cluster counts $(n_1,\ldots,n_m)$.
\end{theorem}

\begin{proof}
We first prove the pointwise estimate
\begin{equation}\label{eq:pointwise}
 |\ind_{I_i}(x)-Q_i(x)|\le C_i x^K P(x)^2\qquad(x>0).
\end{equation}
Outside $I_i$, write $u=x-a_i$, so $|u|\ge r_i$. Away from the other
candidate centers, cancellation gives
\[
 \frac{|Q_i(x)|}{x^KP(x)^2}
 =\frac{a_i^{-K}}{D_i^2}\frac{|1-d_iu|}{u^2}
 \le\frac{a_i^{-K}}{D_i^2}\left(r_i^{-2}+|d_i|r_i^{-1}\right).
\]
At another candidate center both sides of \eqref{eq:pointwise} are zero,
so the same inequality holds there by direct evaluation.

Inside $I_i$, there are no other candidate centers. Define the positive
smooth function $h_i(x)=x^{-K}L_i(x)^{-2}$. At the center,
\[
 h_i(a_i)=a_i^{-K},\qquad h_i'(a_i)=-d_i a_i^{-K}.
\]
Hence $Q_i(x)$ is $x^KL_i(x)^2$ times the first-degree Taylor polynomial
of $h_i$ at $a_i$, while $1=x^KL_i(x)^2h_i(x)$. Since
$P(x)=(x-a_i)D_iL_i(x)$, Taylor's theorem with integral remainder gives
\[
 \frac{|1-Q_i(x)|}{x^KP(x)^2}
 \le\frac1{2D_i^2}\sup_{t\in I_i}|h_i''(t)|
 \quad(x\ne a_i).
\]
Logarithmic differentiation yields the exact expression
\begin{align*}
 \frac{h_i''(t)}{D_i^2}
 &=t^{-K}\prod_{j\ne i}(t-a_j)^{-2}\times{}\\[-2pt]
 &\quad\left\{
 \left(\frac Kt+2\sum_{j\ne i}\frac1{t-a_j}\right)^2
 +\frac K{t^2}+2\sum_{j\ne i}\frac1{(t-a_j)^2}\right\}.
\end{align*}
It is nonnegative and at most $b_i^{-K}A_i(T_i^2+U_i)$, using
$t\ge b_i$ and $|t-a_j|\ge g_{ij}$. At $x=a_i$, both sides of
\eqref{eq:pointwise} vanish. This proves the pointwise bound everywhere.

Absolute convergence permits summation of \eqref{eq:pointwise}, proving
\eqref{eq:countbound}. Every missed occurrence has $x_n\ge\tau$ and
$|x_n-a_i|\ge r_i$ for all $i$, and therefore contributes at least
$\tau^K\prod_i r_i^2$ to $R_K$. Summing gives \eqref{eq:missing}.
A nonnegative integer smaller than one is zero. Finally the true integer
$N_i$ must lie in its certified interval, giving the exact counts.
\end{proof}

\begin{remark}[What a cluster count means]
A count of one proves that the interval contains one simple spectral
occurrence. A count greater than one may represent a repeated value or
several distinct nearby values. No assumption of simple zeros is used.
The no-missing conclusion also excludes spectral points at the interval
endpoints, because the intervals are open.
\end{remark}

\subsection{Certified errors in the input moments}

Write
\[
 x^KP(x)^2=\sum_{j=K}^{K+2m}u_jx^j,
 \qquad Q_i(x)=\sum_{j=K}^{K+2m}v_{ij}x^j,
\]
where the final coefficient of $Q_i$ is zero. Suppose the input satisfies
$|p_j-\widehat p_j|\le\epsilon_j$. Compute
\begin{align}
 \widehat R&=\sum_j u_j\widehat p_j,
 &R^+&=\widehat R+\sum_j|u_j|\epsilon_j,\\
 \widehat S_i&=\sum_jv_{ij}\widehat p_j,
 &E_i&=\sum_j|v_{ij}|\epsilon_j.
\end{align}

\begin{corollary}[Robust finite verification]\label{cor:noise}
Valid input enclosures imply $0\le R_K\le R^+$ and
\[
 |N_i-\widehat S_i|\le E_i+C_iR^+,
 \qquad N_{\rm miss}\le\frac{R^+}{\tau^K\prod_i r_i^2}.
\]
Thus the no-missing test with $R^+$ and a unique integer in each enlarged
count interval constitute a rigorous certificate. One sufficient rounding
test, for a proposed integer $n_i\ge0$, is
\[
 |\widehat S_i-n_i|+E_i+C_iR^+<\frac12.
\]
All verification operations are rational when the centers, radii, threshold,
moment midpoints and error bounds are rational.
\end{corollary}
\begin{proof}
The triangle inequality bounds the errors in both linear moment functionals.
Substitute these bounds into Theorem~\ref{thm:residual}.
\end{proof}

If $R^+<0$, the supplied enclosures contradict the nonnegative residual and
must be rejected. A verifier must not silently replace a negative upper
bound by zero. The input enclosures themselves remain a mathematical
premise: passing this verifier cannot validate arbitrary untrusted moment
data.

\subsection{Why a fixed-width certificate eventually exists}

\begin{theorem}[Eventual certificates]\label{thm:eventual}
Let $\tau>0$ be outside the spectrum. Let $m$ be the number of distinct
spectral values above $\tau$. There is a certificate of the preceding form
using $2m+1$ consecutive moments, with rational candidate centers and radii
in arbitrarily narrow disjoint intervals about those values. If $\tau$ is
rational and the moments can be approximated with certified errors to any
requested precision, an exhaustive rational search eventually finds a
complete correct certificate above $\tau$.
\end{theorem}

\begin{proof}
First let $m\ge1$ and take the candidate centers to be the actual distinct
spectral values above $\tau$. Choose narrow disjoint intervals with
$b_i=a_i-r_i>\tau$. Summability implies that every compact subinterval of
$(0,\infty)$ contains only finitely many spectral occurrences. Since $\tau$
is not a spectral value, all remaining occurrences lie in $(0,\sigma]$ for
some $0<\sigma<\tau$. Write $S_{\rm tail}$ for their first moment.

At the leading spectral values $P$ vanishes exactly. For $0\le x\le\sigma$,
each $a_i>\sigma$, so $P(x)^2\le\prod_i a_i^2$. Therefore
\begin{equation}\label{eq:taildecay}
 R_K\le S_{\rm tail}\,\sigma^{K-1}\prod_i a_i^2.
\end{equation}
For fixed centers and radii, \eqref{eq:Ci} gives
$C_i=O((K+1)^2b_i^{-K})$. Because $\sigma<\tau<b_i$,
\[
 \frac{R_K}{\tau^K\prod_i r_i^2}\longrightarrow0,
 \qquad C_iR_K\longrightarrow0.
\]
The exact integer counts then satisfy the sufficient rounding inequality for
all sufficiently large $K$: indeed
$|S_i-N_i|+C_iR_K\le2C_iR_K\to0$.

Fix one such $K$ with strict slack. The finite formulas are continuous in
the distinct centers, positive admissible radii, and moments. Choose
rational centers and radii sufficiently close to the original ones; their
intervals still enclose the same spectral values. Then choose sufficiently
accurate rational moment enclosures. All strict certificate inequalities
persist. The interval widths may have been arbitrarily small from the start.

When $m=0$, choose $\sigma<\tau$ bounding the entire spectrum. The simple
test $p_K<\tau^K$ eventually holds, since
$p_K\le p_1\sigma^{K-1}$. It certifies that the count above $\tau$ is zero.

Finally enumerate finite rational configurations, positive integers $K$,
and requested moment precisions. For each configuration, the tests of
Corollary~\ref{cor:noise} terminate using rational arithmetic. Dovetailing
this enumeration eventually reaches one of the certificates just proved
to exist. This is an existence and termination argument, not an efficiency
bound.
\end{proof}

\begin{remark}[Conditioning is visible in the theorem]
The proof does not say that fixed approximate centers improve indefinitely
as $K$ grows. Suppose actual leading values $\lambda_i\le B$ have total
multiplicity $M$, $|a_i-\lambda_i|\le\eta$, and the tail is bounded by
$\sigma<\min_i a_i$. A direct product estimate gives
\[
 R_K\le MB^K\eta^2(2B+\eta)^{2m-2}
       +S_{\rm tail}\sigma^{K-1}\prod_i a_i^2.
\]
The first term can require increasingly accurate centers at larger $K$.
Small separation between different centers also enlarges $D_i^{-2}$ and
the $g_{ij}^{-1}$ factors. The explicit constants record this cost.
\end{remark}

\section{The zeta and Dirichlet applications}\label{sec:zeta}

Define
\[
 \xi(s)=\tfrac12s(s-1)\pi^{-s/2}\Gamma(s/2)\zeta(s),
 \qquad \Xi(t)=\xi(\tfrac12+it).
\]
The standard completion satisfies $\xi(s)=\xi(1-s)$ and is entire
\cite{dlmf,cardon}. Use RH, as provided by the inspected development, and list
positive zero ordinates $\gamma_n$ with multiplicity. Set
\begin{equation}\label{eq:zetanodes}
 x_n=\frac1{\tfrac14+\gamma_n^2}.
\end{equation}
This is precisely the positive denominator arising from
$\rho(1-\rho)$ in the repository's transform reduction.

\begin{proposition}[Moment coordinate]\label{prop:xi}
The symmetry quotient
\begin{equation}\label{eq:H}
 H(z)=2\xi\!\left(\frac{1+\sqrt{1+4z}}2\right)
\end{equation}
is a single-valued entire function, despite the square-root notation. It has
the locally uniformly convergent product
\begin{equation}\label{eq:Hproduct}
 H(z)=\prod_n(1+x_nz),\qquad H(0)=1.
\end{equation}
Consequently, for the analytic logarithm at zero,
\begin{equation}\label{eq:logmoments}
 p_j=(-1)^{j-1}j[z^j]\log H(z)
     =\frac{(-1)^{j-1}}{(j-1)!}(\log H)^{(j)}(0),
\end{equation}
and
\begin{equation}\label{eq:p1}
 p_1=1+\frac{\gamma_{\!E}}2-\frac12\log(4\pi).
\end{equation}
\end{proposition}

\begin{proof}
The entire function $w\mapsto\xi(\tfrac12+w)$ is even. Its power series
is therefore an entire power series in $w^2$. Substituting $w^2=z+\tfrac14$
proves that \eqref{eq:H} is entire and independent of the square-root choice.
The paired Hadamard product under RH is
\[
 \xi(\tfrac12+w)=\xi(\tfrac12)
       \prod_{\gamma_n>0}\left(1+\frac{w^2}{\gamma_n^2}\right).
\]
Pairing the genus-one factors removes their exponential factors, while
evenness removes the possible linear exponential prefactor. The usual
order-one factorization gives $\sum_n\gamma_n^{-2}<\infty$. There is no
central zeta zero; for example the alternating eta series at $1/2$ is
positive while $1-\sqrt2<0$, so $\zeta(1/2)\ne0$.

Divide the paired product at $w^2=z+\tfrac14$ by its value at
$w^2=\tfrac14$. Since $\xi(1)=\tfrac12$, this gives
\[
 2\xi(\tfrac12+w)
 =\prod_n\frac{\gamma_n^2+w^2}{\gamma_n^2+\tfrac14},
\]
which is \eqref{eq:Hproduct}. Summability justifies taking the logarithmic
power series near zero, proving \eqref{eq:logmoments}. Finally
$H'(0)/H(0)=\xi'(1)/\xi(1)$. Substitute the Laurent expansion
$\zeta(1+u)=u^{-1}+\gamma_E+O(u)$ into the defining completion, together
with $\Gamma'(1/2)/\Gamma(1/2)=-\gamma_E-2\log2$, to obtain
\eqref{eq:p1}.
\end{proof}

\begin{corollary}[Complete zeta-zero cluster certificates]\label{cor:zeta}
Use certified moment enclosures from \eqref{eq:logmoments}. Suppose
$0<\tau<a_i-r_i<a_i+r_i<4$, and the tests of
Corollary~\ref{cor:noise} succeed. Then every positive-ordinate nontrivial
zeta zero with
\[
 \gamma\le T:=\sqrt{\tau^{-1}-\tfrac14}
\]
lies in one of the disjoint intervals
\begin{equation}\label{eq:gammawindows}
 J_i=\left(
 \sqrt{(a_i+r_i)^{-1}-\tfrac14},\quad
 \sqrt{(a_i-r_i)^{-1}-\tfrac14}\right).
\end{equation}
The certified integer $n_i$ is the number of zeros in $J_i$, counting
multiplicity, and the total count is $\sum_i n_i$. If the threshold is not
a spectral value, arbitrarily narrow such certificates exist.
\end{corollary}
\begin{proof}
The map $x\mapsto\sqrt{x^{-1}-1/4}$ is decreasing on $(0,4)$ and carries
\eqref{eq:zetanodes} back to the positive ordinates. Apply the preceding
theorems. The threshold and boundary exclusions are supplied by the
no-missing certificate.
\end{proof}

For numerical conditioning it is convenient to use $y_n=cx_n$ for a
positive scale $c$. Replace $H(z)$ by $H(cz)$ and replace $x^{-1}$ in
\eqref{eq:gammawindows} by $c/y$. All abstract formulas are unchanged.

\subsection{Scope of the Dirichlet extension}

The abstract theorem applies to any positive trace-class spectrum with
integer multiplicities, using $p_j=\tr(T^j)$. Its Dirichlet application is
also straightforward once the spectral count convention is fixed.
For a primitive nonprincipal character $\chi$, form the product of the
completed $L$-functions for $\chi$ and $\overline\chi$. Their functional
equations make the product even about $s=1/2$. Under Dirichlet GRH its
noncentral zeros lie on that line. Remove its central factor first, and
apply the even paired canonical product as above. The resulting positive
moments certify combined noncentral zero counts for the pair.

For a real character, using the character twice doubles those counts; one
can instead use the appropriately normalized single completion after
removing its central zero. The zeta completion handles the primitive
principal character. Imprimitive characters should be reduced to their
primitive inducing characters, treating their additional Euler factors
separately. Central vanishing is not excluded by Dirichlet GRH and must
not be silently discarded. We do not claim that the inspected repository
proves GRH for all automorphic or Artin $L$-functions.

\subsection{Relation to earlier reconstruction results}

Finite moment reconstruction and polynomial dual certificates are
established methods. De Castro and Gamboa prove recovery of a nonnegative
measure on $s$ points from $2s+1$ generalized moments \cite{decg}.
Batenkov and Yomdin quantify the conditioning of finite Prony systems
\cite{prony}; Cand\`es and Fernandez-Granda give stable recovery results
from noisy bandlimited data under separation hypotheses \cite{superres}.
Cardon and de Gaston connect orthogonal polynomials to RH with simple
zeros \cite{cardon}.

A particularly close recent comparison is \emph{A finite arithmetic test
for simplicity of the Riemann zeros, after Suzuki} \cite{ald}. It compares
multiplicity-weighted Taylor moments with finite operator-compression
traces to detect excess multiplicity. The present certificate uses direct
positive moments, a candidate polynomial, and scalar Hermite bounds to
certify complete cluster counts. It makes no distinction between a repeated
root and multiple roots within a cluster of count greater than one.

The specific contribution proposed here is the combination of an infinite
summable integer spectrum, an explicit residual and count constant, a
fixed-width shifted moment window, rigorous input-error propagation, and
eventual rational certificates. None of the cited classical ingredients is
claimed as new, and the literature search does not establish priority for
the combined formulation.
